\subsection{Формулы конечных приращений Лагранжа и Коши.}

\subsubsection*{Теорема Лагранжа (о конечном приращении)}

\begin{tcolorbox}[title=Теорема L (Лагранжа)]
    Пусть функция $f(x)$:
    \begin{enumerate}
        \item Непрерывна на отрезке $[a, b]$ ($f(x) \in C[a, b]$).
        \item Дифференцируема на интервале $(a, b)$.
    \end{enumerate}
    Тогда существует точка $\xi \in (a, b)$, такая что:
    \[ \frac{f(b) - f(a)}{b - a} = f'(\xi) \]
    Или в виде формулы конечных приращений:
    \[ f(b) - f(a) = f'(\xi)(b - a) \]
\end{tcolorbox}

\textbf{Доказательство:}
Введем вспомогательную функцию:
\[ F(x) = f(x) - \frac{f(b)-f(a)}{b-a}(x-a) \]
Проверим значения на концах отрезка:
1) $F(a) = f(a) - \frac{f(b)-f(a)}{b-a} \cdot 0 = f(a)$.
2) $F(b) = f(b) - \frac{f(b)-f(a)}{b-a}(b-a) = f(b) - (f(b) - f(a)) = f(a)$.

Так как $F(a) = F(b)$, функция $F(x)$ удовлетворяет всем условиям \textbf{теоремы Ролля}.
Следовательно, $\exists \xi \in (a, b)$, такая что $F'(\xi) = 0$.

Найдем производную $F'(x)$:
\[ F'(x) = f'(x) - \frac{f(b)-f(a)}{b-a} \cdot 1 \]
Подставим точку $\xi$:
\[ F'(\xi) = f'(\xi) - \frac{f(b)-f(a)}{b-a} = 0 \]
Откуда получаем утверждение теоремы:
\[ \frac{f(b)-f(a)}{b-a} = f'(\xi). \]

\textbf{Оценка приращения:}
\[ |f(b) - f(a)| \le \max_{[a, b]} |f'(\xi)| \cdot |b - a| \]

\textbf{Следствия из теоремы Лагранжа:}
\begin{enumerate}
    \item Если $f'(x) \ge 0$ на $(a, b)$, то $f(x)$ возрастает ($f(x) \nearrow$).
    Если $f'(x) \le 0$ на $(a, b)$, то $f(x)$ убывает ($f(x) \searrow$).
    \item Если $f'(x) \equiv 0$ на $(a, b)$, то $f(x) \equiv \text{const}$ на $[a, b]$.
\end{enumerate}

\subsubsection*{Теорема Коши (обобщенная формула конечных приращений)}

\begin{tcolorbox}[title=Теорема C (Коши)]
    Пусть функции $x = x(t)$ и $y = y(t)$:
    \begin{enumerate}
        \item Непрерывны на $[a, b]$ ($x, y \in C[a, b]$).
        \item Дифференцируемы на $(a, b)$.
        \item $x'(t) \neq 0$ на $(a, b)$.
    \end{enumerate}
    Тогда существует точка $\xi \in (a, b)$, такая что:
    \[ \frac{y(b) - y(a)}{x(b) - x(a)} = \frac{y'(\xi)}{x'(\xi)} \]
    Или в симметричной форме:
    \[ y'(\xi)(x(b) - x(a)) = x'(\xi)(y(b) - y(a)) \]
\end{tcolorbox}

\textit{Замечание:} Теорема Ролля напрямую к траектории на плоскости не применима (возможны петли, где производные не обнуляются одновременно), поэтому требуется доказательство через вспомогательную функцию.

\textbf{Доказательство:}
Рассмотрим вспомогательную функцию $F(t)$:
\[ F(t) = x(t)(y(b) - y(a)) - y(t)(x(b) - x(a)) \]
Проверим значения на концах отрезка:
\begin{itemize}
    \item $F(a) = x(a)(y(b)-y(a)) - y(a)(x(b)-x(a)) = x(a)y(b) - \underline{x(a)y(a)} - y(a)x(b) + \underline{y(a)x(a)} = x(a)y(b) - y(a)x(b)$.
    \item $F(b) = x(b)(y(b)-y(a)) - y(b)(x(b)-x(a)) = \underline{x(b)y(b)} - x(b)y(a) - \underline{y(b)x(b)} + y(b)x(a) = y(b)x(a) - x(b)y(a)$.
\end{itemize}
Видим, что $F(a) = F(b)$.
По теореме Ролля $\exists \xi \in (a, b)$, такая что $F'(\xi) = 0$.
\[ F'(\xi) = x'(\xi)(y(b) - y(a)) - y'(\xi)(x(b) - x(a)) = 0 \]
Откуда:
\[ \frac{y(b) - y(a)}{x(b) - x(a)} = \frac{y'(\xi)}{x'(\xi)} \]
Теорема доказана.

\textit{Замечание:} Теорема Лагранжа — это частный случай теоремы Коши (при $x(t) = t$).

\subsubsection*{Обобщение для производных высших порядков}

\begin{tcolorbox}[title=Утверждение (Многократное применение теоремы Коши)]
    Пусть $\varphi(t), \psi(t) \in C[a, b]$ и $n$ раз дифференцируемы на $(a, b)$.
    Пусть выполнены условия зануления производных в точке $a$:
    \[ \varphi(a) = \varphi'(a) = \dots = \varphi^{(n-1)}(a) = 0 \]
    \[ \psi(a) = \psi'(a) = \dots = \psi^{(n-1)}(a) = 0 \]
    И пусть $\psi^{(k)}(t) \neq 0$ на $(a, b)$.
    
    Тогда существует $\xi \in (a, b)$, такая что:
    \[ \frac{\varphi(b)}{\psi(b)} = \frac{\varphi^{(n)}(\xi)}{\psi^{(n)}(\xi)} \]
\end{tcolorbox}

\textbf{Доказательство:}
Применяем теорему Коши последовательно.
Так как $\psi(b) \neq 0$ (иначе по т. Ролля $\psi'(\xi)=0$, противоречие с условием), запишем отношение:
\[ \frac{\varphi(b)}{\psi(b)} = \frac{\varphi(b) - \varphi(a)}{\psi(b) - \psi(a)} = \frac{\varphi'(\xi_1)}{\psi'(\xi_1)}, \quad \text{где } \xi_1 \in (a, b). \]
Теперь применим теорему Коши к отрезку $[a, \xi_1]$:
\[ \frac{\varphi'(\xi_1)}{\psi'(\xi_1)} = \frac{\varphi'(\xi_1) - \varphi'(a)}{\psi'(\xi_1) - \psi'(a)} = \frac{\varphi''(\xi_2)}{\psi''(\xi_2)}, \quad \text{где } \xi_2 \in (a, \xi_1). \]
Продолжая этот процесс $n$ раз, получаем цепочку равенств:
\[ \dots = \frac{\varphi^{(n-1)}(\xi_{n-1})}{\psi^{(n-1)}(\xi_{n-1})} = \frac{\varphi^{(n-1)}(\xi_{n-1}) - \varphi^{(n-1)}(a)}{\psi^{(n-1)}(\xi_{n-1}) - \psi^{(n-1)}(a)} = \frac{\varphi^{(n)}(\xi)}{\psi^{(n)}(\xi)}, \]
где $\xi \in (a, \xi_{n-1}) \subset \dots \subset (a, b)$.
Утверждение доказано.
